?
\begin{tcolorbox}[colback=red!3!white,colframe=red!50!black,boxrule=0.8pt,title={Correction notes},fonttitle=\bfseries]
\textbf{Corrections from the calculation check (30~Sept.~2026).} The following places are corrected or annotated in the text; each carries a dated note with the previous value:
\begin{itemize}
\item Doc.~025: $\frac{16}{9}\xi^2 \cdot 7500 = 2.35\times10^{-4} \to 2.37\times10^{-4}$; $\tau_{\text{Molecule}} = \tau_\xi/\xi^3$ $10^2~\text{s} \to 0.14$~s, $\tau_{\text{Cell}} = \tau_\xi/\xi^4$ $10^9~\text{s} \approx 30$~years $\to 1.1\times10^3~\text{s} \approx 18$~min; note on the dimension of $z = \xi \cdot d$.
\item Doc.~025 (1~Oct.~2026): $\rhoCMB = 4.87 \times 10^{41} \to 2.0 \times 10^{-15}$~eV$^4$ (section on the $\xi$ length scale), as in Doc.~061.
\item Doc.~025 (1~Oct.~2026): Casimir energy density $\pi^2\hbar c/(240\,d^4) \to \pi^2\hbar c/(720\,d^4)$ ($240$ belongs to the pressure), ratio to $\rhoCMB$ $308 \to 102.8$ (text, table, prediction); ``Experimental value 312, agreement 98.7\,\%'' $\to$ $103.8$ with rounded $\Lxi$, not an independent measured value.
\end{itemize}
\end{tcolorbox}

\section*{Abstract}
		This document presents the cosmological aspects of the T0-Theory with the universal $\xi$-parameter as the foundation for a static, eternally existing universe. Starting from the time-energy duality, it is argued that a classical Big Bang singularity is excluded and replaced by a tiny but finite core with minimal length scale $L_0$ (from $\xi$), and it is proposed that the cosmic microwave background radiation (CMB) as well as the Casimir effect be understood as two manifestations of the same $\xi$-field. As the sixth document of the T0 series, it carries the basic principles of the preceding documents over to cosmological questions.

	\medskip\noindent\textbf{Note on versions.} This English version of the document is more extensive than the German one. Earlier revisions shortened or changed parts of the German version. Where the two differ, this more extensive English version applies.


\medskip\noindent\textbf{Status markers.} Statements marked \textbf{[S]} are \emph{speculative}: a hint or conjecture that has not yet been derived or numerically checked in the corpus. The present corpus deliberately sets the cosmic sector ($H_0$, $\Lambda$, dark energy, origin of the CMB) aside, because the Standard Model does not derive it either (P39); the cosmological interpretations in this document should be read in that sense.

	
	
	\section{Introduction}
	
	\subsection{Cosmology within the Framework of the T0-Theory}
	
	The T0-Theory introduces a relationship between the microscopic quantum vacuum and macroscopic cosmic structures. The cosmological quantities treated in this document are traced back to the parameter $\xipar = \frac{4}{3} \times 10^{-4}$.
	
	\begin{keyresult}
		\textbf{Central Thesis of T0-Cosmology [S]:}
		
		The universe is static and eternally existing. The observed cosmic phenomena are interpreted here as manifestations of the $\xi$-field, not as a consequence of spacetime expansion (cosmic sector, cf.\ P39).
	\end{keyresult}
	
	\subsection{Connection to the T0 Document Series}
	
	This cosmological analysis builds on the fundamental insights of the previous T0 documents:
	
	\begin{itemize}
		\item \textbf{003\_T0\_Grundlagen\_v1\_En.pdf:} Geometric parameter $\xipar$ and fractal spacetime structure
		\item \textbf{011\_T0\_Feinstruktur\_En.pdf:} Electromagnetic interactions in the $\xi$-field
		\item \textbf{012\_T0\_Gravitationskonstante\_En.pdf:} Gravitation theory from $\xi$-geometry
		\item \textbf{006\_T0\_Teilchenmassen\_En.pdf:} Mass spectrum as the basis for cosmic structure formation
		\item \textbf{007\_T0\_Neutrinos\_En.pdf:} Neutrino oscillations in cosmic dimensions
	\end{itemize}
	
	\section{Time-Energy Duality and the Static Universe}
	
	\subsection{Heisenberg's Uncertainty Principle as a Cosmological Principle}
	
	\begin{revolutionary}[Key statement]
		\textbf{Starting argument of this document:}
		
		According to the argument presented here, Heisenberg's uncertainty principle $\Delta E \times \Delta t \geq \frac{\hbar}{2}$ excludes a classical Big Bang singularity with infinite density; its place is taken by a tiny but finite core with minimal length scale $L_0$ (from $\xi$). The corpus keeps this conclusion; the justification via the uncertainty principle, however, has been replaced by the corpus's own time-energy reciprocity $T\cdot E = 1$ (Doc.~306).
	\end{revolutionary}
	
	In natural units ($\hbar = c = k_B = 1$), the time-energy uncertainty relation reads:
	
	\begin{equation}
		\Delta E \times \Delta t \geq \frac{1}{2}
	\end{equation}
	
	In the original version of this document, the chain of argument reads:
	
	\begin{itemize}
		\item A temporal beginning (Big Bang) would imply $\Delta t$ = finite
		\item This leads to $\Delta E \to \infty$ - physically inconsistent
		\item Therefore, the universe must have existed eternally: $\Delta t = \infty$
		\item The universe is static, without expanding space
	\end{itemize}
	
	\subsection{Consequences for Standard Cosmology}
	
	\begin{warning}
		\textbf{Problems of Big Bang Cosmology as Seen in This Document:}
		
		\begin{enumerate}
			\item \textbf{Conflict with Quantum Mechanics:} Finite $\Delta t$ requires infinite energy (argument of this document; justification replaced by Doc.~306)
			\item \textbf{Fine-Tuning Problems:} Over 20 free parameters required
			\item \textbf{Dark Matter/Energy:} 95\% unknown components
			\item \textbf{Hubble Tension:} 9\% discrepancy between local and cosmic measurements
			\item \textbf{Age Problem:} Objects older than the stated age of the universe
		\end{enumerate}
	\end{warning}
	
	\section{The Cosmic Microwave Background Radiation (CMB)}
	
	\subsection{CMB as $\xi$-Field Manifestation}
	
	\textbf{[S]} If there is no Big Bang, the CMB needs an origin other than the z=1100 decoupling of standard cosmology. The T0-Theory proposes to explain the CMB through $\xi$-field quantum fluctuations (cosmic sector, cf.\ P39).
	
	\begin{formula}
		\textbf{T0-CMB-Temperature Relation:}
		\begin{equation}
			\frac{T_{\text{CMB}}}{\Exi} = \frac{16}{9} \xipar^2
		\end{equation}
	\end{formula}
	
	With $\Exi = \frac{1}{\xipar} = \frac{3}{4} \times 10^4$ (natural units) and $\xipar = \frac{4}{3} \times 10^{-4}$, the result is:
	
	\begin{align}
		T_{\text{CMB}} &= \frac{16}{9} \xipar^2 \times \Exi \\
		&= \frac{16}{9} \times \left(\frac{4}{3} \times 10^{-4}\right)^2 \times \frac{3}{4} \times 10^4 \\
		&= \frac{16}{9} \times 1.78 \times 10^{-8} \times 7500 \\
		&= 2.37 \times 10^{-4} \text{ (natural units)}
	\end{align}
	\textit{(Corrected on 30~Sept.~2026: previously $2.35 \times 10^{-4}$ -- $\frac{16}{9} \cdot \left(\frac{4}{3}\cdot10^{-4}\right)^2 \cdot 7500 = 2.370\times10^{-4}$, with the rounded $1.78\times10^{-8}$: $2.373\times10^{-4}$; this would correspond to $2.75$~K.)}
	
	\textbf{Conversion to SI Units:} $T_{\text{CMB}} = 2.725$ K
	
	This is close to the Planck measurement; the dimensionless ratios differ by about 1\,\% ($3.13$ versus $3.16 \times 10^{-8}$, see the table of dimensionless $\xi$-ratios).
	
	\subsection{CMB Energy Density and Characteristic Length Scale}
	
	The CMB energy density defines a fundamental characteristic length scale of the $\xi$-field:
	
	\begin{equation}
		\rhoCMB = \frac{\xipar}{\Lxi^4}
	\end{equation}
	
	From this follows the characteristic $\xi$-length scale:
	
	\begin{equation}
		\Lxi = \left(\frac{\xipar}{\rhoCMB}\right)^{1/4}
	\end{equation}
	
	\begin{keyresult}
		\textbf{Characteristic $\xi$-Length Scale:}
		
		Using the experimental CMB data, the result is:
		\begin{equation}
			\Lxi = 100 \, \mu\text{m}
		\end{equation}
		
		This length scale marks the transition region between microscopic quantum effects and macroscopic cosmic phenomena.
	\end{keyresult}
	
	\section{Casimir Effect and $\xi$-Field Connection}
	
	\subsection{Casimir-CMB Ratio as Experimental Comparison}
	
	The ratio between Casimir energy density and CMB energy density is compatible with the characteristic $\xi$-length scale and supports interpreting both quantities through the same $\xi$-field.
	
	The Casimir energy density at plate separation $d = \Lxi$ is:
	
	\begin{equation}
		|\rhoCasimir| = \frac{\pi^2 \hbar c}{720 \times \Lxi^4}
	\end{equation}
	
	The theoretical ratio yields:
	
	\begin{equation}
		\frac{|\rhoCasimir|}{\rhoCMB} = \frac{\pi^2}{720 \xipar} = \frac{\pi^2 \times 10^4}{960} \approx 102.8
	\end{equation}
	
	\textit{(Corrected on 1~Oct.~2026: previously $\pi^2\hbar c/(240\,\Lxi^4)$ and ratio $\pi^2/(240\,\xi) = \pi^2 \times 10^4/320 \approx 308$ -- $\pi^2\hbar c/(240\,d^4)$ is the Casimir pressure; the energy density between the plates is $\pi^2\hbar c/(720\,d^4)$, hence $\pi^2/(720\,\xi) \approx 102.8$. Since $\Lxi$ is fixed from $\rhoCMB$ via $\rhoCMB = \xi\hbar c/\Lxi^4$, the ratio at $d = \Lxi$ is an identity and not an independent comparison.)}
	
	\begin{experiment}
		\textbf{Numerical Comparison:}
		
		The Python verification script \texttt{CMB\_En.py} (available on GitHub: ) yields:
		
		\begin{itemize}
			\item Theoretical value: 102.8
			\item The same expression in SI units with rounded $\Lxi = 100\,\mu$m: 103.8
			\item Deviation: 1.0\,\% (rounding of $\Lxi$)
		\end{itemize}
		\textit{(Corrected on 1~Oct.~2026: previously ``Theoretical Prediction: 308, Experimental Value: 312, Agreement: 98.7\,\%'' -- 312 is not a measured value but the same formula with $\Lxi = 100\,\mu$m instead of $100.24\,\mu$m (from $\rhoCMB = 4.175 \times 10^{-14}$~J/m$^3$); with the energy density ($720$ instead of $240$) one obtains $102.8$ and $103.8$.)}
	\end{experiment}
	
	\subsection{$\xi$-Field as Universal Vacuum}
	
	\begin{revolutionary}[Key statement]
		\textbf{Interpretation [S]:}
		
		According to this interpretation, the $\xi$-field manifests itself both in the free CMB radiation and in the geometrically confined Casimir vacuum. The document reads this as a hint that the $\xi$-field acts as the universal quantum vacuum; the numerical comparison is an indication, not a proof.
	\end{revolutionary}
	
	The characteristic $\xi$-length scale $\Lxi$ is the point where CMB vacuum energy density and Casimir energy density reach comparable orders of magnitude:
	
	\begin{align}
		\text{Free Vacuum:} \quad &\rhoCMB = \frac{\pi^2}{15}\,T_{\text{CMB}}^4 = 2.0 \times 10^{-15}\,\text{eV}^4 \text{ (natural units)} \\
		\text{Confined Vacuum:} \quad &|\rhoCasimir| = \frac{\pi^2}{720 d^4}
	\end{align}
	\textit{(Corrected on 1~Oct.~2026: previously $\pi^2/(240\,d^4)$ -- energy density instead of pressure, see above. Likewise previously $\rhoCMB = +4.87 \times 10^{41}$ -- this value cannot be reproduced in any unit system; with $T_{\text{CMB}} = 2.35 \times 10^{-4}$~eV one has $\rhoCMB = (\pi^2/15)\,T_{\text{CMB}}^4 = 2.0 \times 10^{-15}$~eV$^4$, as corrected in Doc.~061.)}
	
	\section{Cosmic Redshift: Alternative Interpretations}
	
	\subsection{The Mathematical Model of the T0-Theory}
	
	The T0-Theory provides a mathematical model for the observed cosmic redshift that \emph{allows alternative interpretations}, without committing to a specific physical cause.
	
	\begin{formula}
		\textbf{Fundamental T0-Redshift Model:}
		\begin{equation}
			z = \xi \cdot d \quad \text{(see Document~026)}
		\end{equation}
		where $\lambda_0$ is the emitted wavelength, $d$ the distance, and $\Exi$ the characteristic $\xi$-energy. \textit{(Note of 30~Sept.~2026: $\xi$ is dimensionless, $d$ a length; $z = \xi \cdot d$ is not dimensionless without a length scale, not even for $\hbar = c = 1$. Doc.~026 writes $z \approx d \cdot C \cdot \xi$; the constant $C$ is missing here. $\lambda_0$ and $\Exi$ do not appear in the formula.)}
	\end{formula}
	
	\subsection{Alternative Physical Interpretations}
	
	The same mathematical model can be realized through different physical mechanisms:
	
	\begin{alternative}
		\textbf{Interpretation 1: Energy Loss Mechanism}
		
		Photons lose energy through interaction with the omnipresent $\xi$-field:
		\begin{equation}
			\frac{dE}{dx} = -\frac{\xipar E^2}{\Exi}
		\end{equation}
		
		\textbf{Physical Assumptions:}
		\begin{itemize}
			\item Direct energy transfer from the photon to the $\xi$-field
			\item Continuous process over cosmic distances
			\item No space expansion required
		\end{itemize}
	\end{alternative}
	
	\begin{alternative}
		\textbf{Interpretation 2: Gravitational Deflection by Mass}
		
		The redshift arises from cumulative gravitational deflection effects along the light path:
		\begin{equation}
			z(d) = \int_0^d \xi \cdot \rho_{\text{Matter}}(x)\, dx \quad \text{(see Document~026)}
		\end{equation}
		
		\textbf{Physical Assumptions:}
		\begin{itemize}
			\item Matter distribution determined by $\xi$-parameter
			\item Gravitational frequency shift accumulates over distance
			\item Static universe with homogeneous matter distribution
		\end{itemize}
	\end{alternative}
	
	\begin{alternative}
		\textbf{Interpretation 3: Spacetime Geometry Effects}
		
		The $\xi$-field structure of spacetime modifies light propagation:
		\begin{equation}
			ds^2 = \left(1 + \frac{\xipar \lambda_0}{\Exi}\right) dt^2 - dx^2
		\end{equation}
		
		\textbf{Physical Assumptions:}
		\begin{itemize}
			\item Wavelength-dependent metric coefficients
			\item $\xi$-field as fundamental spacetime component
			\item Geometric cause of frequency shift
		\end{itemize}
	\end{alternative}
	
	\subsection{Experimental Distinction of Interpretations}
	
	\begin{experiment}
		\textbf{Tests to Distinguish Mechanisms:}
		
		\begin{enumerate}
			\item \textbf{Polarization Analysis:}
			\begin{itemize}
				\item Energy Loss: No polarization effects
				\item Gravitational Deflection: Weak polarization rotation
				\item Geometric Effects: Specific polarization patterns
			\end{itemize}
			
			\item \textbf{Temporal Variation:}
			\begin{itemize}
				\item Energy Loss: Constant effect
				\item Gravitational Deflection: Varies with local matter density
				\item Geometric Effects: Dependent on $\xi$-field fluctuations
			\end{itemize}
			
			\item \textbf{Spectral Signatures:}
			\begin{itemize}
				\item Energy Loss: Smooth wavelength-dependent curve
				\item Gravitational Deflection: Discrete peaks at mass concentrations
				\item Geometric Effects: Interference patterns at characteristic frequencies
			\end{itemize}
		\end{enumerate}
	\end{experiment}
	
	\subsection{Common Predictions of All Interpretations}
	
	Regardless of the specific mechanism, the T0 model predicts:
	
	\begin{keyresult}
		\textbf{Universal T0-Redshift Predictions [S] (original state; cf.\ K6 in Doc.~190):}
		
		\begin{itemize}
			\item \textbf{Wavelength Dependence:} $z \propto \lambda_0$
			\item \textbf{Distance Dependence:} $z \propto d$ (linear, not exponential)
			\item \textbf{Characteristic Scale:} Effects maximal at $\lambda \sim \Lxi$
			\item \textbf{Ratio of Different Wavelengths:} $z_1/z_2 = \lambda_1/\lambda_2$
		\end{itemize}
	\end{keyresult}
	
	\subsection{Strategic Significance of Multiple Interpretations}
	
	\begin{warning}
		\textbf{Methodological Advantage:}
		
		By offering multiple interpretations, the T0-Theory avoids:
		\begin{itemize}
			\item Premature commitment to a specific mechanism
			\item Exclusion of experimentally equivalent explanations
			\item Ideological preferences over physical evidence
			\item Limitation of future theoretical developments
		\end{itemize}
		
		This corresponds to the principle of scientific objectivity and falsifiability.
	\end{warning}	
	\section{Structure Formation in the Static $\xi$-Universe}
	
	\subsection{Continuous Structure Development}
	
	In the static T0-universe, structure formation occurs continuously without Big Bang constraints:
	
	\begin{equation}
		\frac{d\rho}{dt} = -\nabla \cdot (\rho \mathbf{v}) + S_\xi(\rho, T, \xipar)
	\end{equation}
	
	where $S_\xi$ is the $\xi$-field source term for continuous matter/energy transformation.
	
	\subsection{$\xi$-Supported Continuous Creation}
	
	The $\xi$-field enables continuous matter/energy transformation:
	
	\begin{align}
		\text{Quantum Vacuum} &\xrightarrow{\xipar} \text{Virtual Particles} \\
		\text{Virtual Particles} &\xrightarrow{\xipar^2} \text{Real Particles} \\
		\text{Real Particles} &\xrightarrow{\xipar^3} \text{Atomic Nuclei} \\
		\text{Atomic Nuclei} &\xrightarrow{\text{Time}} \text{Stars, Galaxies}
	\end{align}
	
	The energy balance is maintained by:
	
	\begin{equation}
		\rho_{\text{total}} = \rho_{\text{Matter}} + \rho_{\xi\text{-Field}} = \text{constant}
	\end{equation}
	
	\subsection{Solution to Structure Formation Problems}
	
	\begin{keyresult}
		\textbf{Advantages of T0 Structure Formation [S]:}
		
		\begin{itemize}
			\item \textbf{Unlimited Time:} Structures can become arbitrarily old
			\item \textbf{No Fine-Tuning:} Continuous evolution instead of critical initial conditions
			\item \textbf{Hierarchical Development:} From quantum fluctuations to galaxy clusters
			\item \textbf{Stability:} Static universe prevents cosmic catastrophes
		\end{itemize}
	\end{keyresult}
	
	\section{Dimensionless $\xi$-Hierarchy}
	
	\subsection{Energy Scale Ratios}
	
	The $\xi$-relations can be written as dimensionless ratios:
	
	\begin{longtable}{lcc}
		\caption{Dimensionless $\xi$-Ratios in Cosmology} \\
		\toprule
		\textbf{Ratio} & \textbf{Expression} & \textbf{Value} \\
		\midrule
		\endfirsthead
		\multicolumn{3}{c}{\tablename\ \thetable{} -- Continued} \\
		\toprule
		\textbf{Ratio} & \textbf{Expression} & \textbf{Value} \\
		\midrule
		\endhead
		CMB Temperature & $\frac{T_{\text{CMB}}}{\Exi}$ & $3.13 \times 10^{-8}$ \\
		Theory & $\frac{16}{9}\xipar^2$ & $3.16 \times 10^{-8}$ \\
		Characteristic Length & $\frac{\ell_{\xipar}}{\Lxi}$ & $\xipar^{-1/4}$ \\
		Casimir-CMB & $\frac{|\rhoCasimir|}{\rhoCMB}$ & $\frac{\pi^2 \times 10^4}{960}$ \\
		Hubble Substitute & $\frac{\xipar x}{\Exi \lambda}$ & dimensionless \\
		Structure Scale & $\frac{L_{\text{Structure}}}{\Lxi}$ & $(\text{Age}/\tau_\xi)^{1/4}$ \\
		\bottomrule
	\end{longtable}
	\textit{(Corrected on 1~Oct.~2026: row Casimir-CMB previously $\pi^2 \times 10^4/320$ -- with the energy density $\pi^2/(720\,\xi) = \pi^2 \times 10^4/960 \approx 102.8$.)}
\normalsize
	
	\begin{warning}
		\textbf{Structure of the T0-Cosmology Relations:}
		
		The theoretical $\xi$-relations contain only simple numbers:
		\begin{itemize}
			\item Fractions: $\frac{4}{3}$, $\frac{3}{4}$, $\frac{16}{9}$
			\item Powers of Ten: $10^{-4}$, $10^3$, $10^4$
			\item Mathematical Constants: $\pi^2$
		\end{itemize}
		
		No fitted decimal numbers appear on the theory side; the expressions follow from $\xi$ and integer structural choices that enter as motivated settings.
	\end{warning}
	
	\section{Experimental Predictions and Tests}
	
	\subsection{Precision Casimir Measurements}
	
	\begin{experiment}
		\textbf{Critical Test at Characteristic Length Scale:}
		
		Casimir force measurements at $d = 100\,\mu$m should show the theoretical ratio 102.8:1 to the CMB energy density. \textit{(Corrected on 1~Oct.~2026: previously $308{:}1$ -- that is the ratio of the Casimir pressure $\pi^2\hbar c/(240\,d^4)$ to $\rhoCMB$; for the energy density $\pi^2/(720\,\xi) \approx 102.8$ holds. Both numbers follow from the definition of $\Lxi$ and the standard Casimir formula; a deviation would test the Casimir formula, not the $\xi$-field.)}
		
		\textbf{Experimental Accessibility:} $\Lxi = 100\,\mu$m is within the measurable range of modern Casimir experiments.
	\end{experiment}
	
	\subsection{Electromagnetic $\xi$-Resonance}
	
	Maximum $\xi$-field-photon coupling at characteristic frequency:
	
	\begin{equation}
		\nu_\xi = \frac{c}{\Lxi} = \frac{3 \times 10^8}{10^{-4}} = 3 \times 10^{12} \text{ Hz} = 3 \text{ THz}
	\end{equation}
	
	At this frequency, electromagnetic anomalies should occur, measurable with high-precision THz spectrometers.
	
	\subsection{Cosmic Tests of Wavelength-Dependent Redshift}
	
	\textbf{Note:} According to the present state of the corpus, the redshift is achromatic (Doc.~190, K6). The wavelength-dependent tests below reflect the original state of this document.
	
	\begin{experiment}
		\textbf{Multi-Wavelength Astronomy:}
		
		\begin{enumerate}
			\item \textbf{Galaxy Spectra:} Comparison of UV, optical, and radio redshifts
			\item \textbf{Quasar Observations:} Wavelength dependence at high z values
			\item \textbf{Gamma-Ray Bursts:} Extreme UV redshift vs. radio components
		\end{enumerate}
		
		The T0-Theory predicts specific ratios that deviate from standard cosmology.
	\end{experiment}
	
	\section{Cosmological Problems Compared}
	
	\subsection{Comparison: $\Lambda$CDM vs. T0 Model}
	
	The right-hand column gives the interpretation in the static T0 picture; like the whole cosmic sector, it is speculative \textbf{[S]} (P39).
	
	\begin{longtable}{p{3.2cm}p{3.6cm}p{3.6cm}}
		\caption{Cosmological Problems: Standard vs. T0} \\
		\toprule
		\textbf{Problem} & \textbf{$\Lambda$CDM} & \textbf{T0 Interpretation} \\
		\midrule
		\endfirsthead
		\multicolumn{3}{c}{\tablename\ \thetable{} -- Continued} \\
		\toprule
		\textbf{Problem} & \textbf{$\Lambda$CDM} & \textbf{T0 Interpretation} \\
		\midrule
		\endhead
		Horizon Problem & Inflation required & Infinite causal connectivity \\
		Flatness Problem & Fine-tuning & Geometry stabilized over infinite time \\
		Monopole Problem & Topological defects & Defects dissipate over infinite time \\
		Lithium Problem & Nucleosynthesis discrepancy & Nucleosynthesis over unlimited time \\
		Age Problem & Objects older than universe & Objects can be arbitrarily old \\
		$H_0$ Tension & 9\% discrepancy & No $H_0$ in static universe \\
		Dark Energy & 69\% of energy density & Not required \\
		Dark Matter & 26\% of energy density & $\xi$-field effects \\
		\bottomrule
	\end{longtable}
\normalsize
	
	\subsection{Parameter Count}
	
	\begin{revolutionary}[Key statement]
		\textbf{From 25+ Parameters to a Single One:}
		
		\begin{itemize}
			\item Standard Model of Particle Physics: 19+ parameters
			\item $\Lambda$CDM Cosmology: 6 parameters
			\item \textbf{T0-Theory: 1 Parameter ($\xipar$)}, plus integer structural choices as motivated settings and SI anchors for unit conversion
		\end{itemize}
		
		This corresponds to a reduction of the continuous parameters by 96\%. According to the present state of the corpus, however, the cosmological sector requires its own, independently set scale (P39, Doc.~306).
	\end{revolutionary}
	
	\section{Cosmic Timescales and $\xi$-Evolution}
	
	\subsection{Characteristic Timescales}
	
	The $\xi$-field defines fundamental timescales for cosmic processes:
	
	\begin{equation}
		\tau_\xi = \frac{\Lxi}{c} = \frac{10^{-4}}{3 \times 10^8} = 3.3 \times 10^{-13} \text{ s}
	\end{equation}
	
	Longer timescales arise from $\xi$-hierarchies:
	
	\begin{align}
		\tau_{\text{Atom}} &= \frac{\tau_\xi}{\xipar^2} \approx 10^{-5} \text{ s} \\
		\tau_{\text{Molecule}} &= \frac{\tau_\xi}{\xipar^3} \approx 0.14 \text{ s} \\
		\tau_{\text{Cell}} &= \frac{\tau_\xi}{\xipar^4} \approx 1.1 \times 10^3 \text{ s} \approx 18 \text{ min}
	\end{align}
	\textit{(Corrected on 30~Sept.~2026: previously $\tau_{\text{Molecule}} \approx 10^2$~s and $\tau_{\text{Cell}} \approx 10^9~\text{s} \approx 30$~years -- with $\tau_\xi = 3.34\times10^{-13}$~s one has $\tau_\xi/\xi^2 = 1.9\times10^{-5}$~s, $\tau_\xi/\xi^3 = 0.14$~s, $\tau_\xi/\xi^4 = 1.06\times10^3~\text{s} = 17.6$~min.)}
	
	\subsection{Cosmic $\xi$-Cycles}
	
	In this picture, the static T0-universe undergoes $\xi$-driven cycles \textbf{[S]}:
	
	\begin{enumerate}
		\item \textbf{Matter Accumulation:} $\xi$-field → particles → structures
		\item \textbf{Structure Maturity:} Galaxies, stars, planets
		\item \textbf{Energy Return:} Hawking radiation → $\xi$-field
		\item \textbf{Cycle Restart:} New matter generation
	\end{enumerate}
	
	\section{Connection to Dark Matter and Dark Energy}
	
	\subsection{$\xi$-Field as Dark Matter Alternative}
	
	\begin{keyresult}
		\textbf{$\xi$-Field as an Interpretation of Dark Matter [S]:}
		
		\begin{itemize}
			\item Gravitationally acting through energy-momentum tensor
			\item Electromagnetically neutral (detectable only via specific resonances)
			\item Correct cosmological energy density at $\Delta m \sim \xipar \times m_{\text{Planck}}$
			\item Intended to explain galaxy rotation curves without new particles (not worked out quantitatively here)
		\end{itemize}
	\end{keyresult}
	
	\subsection{No Dark Energy Required}
	
	In the static T0-universe, no dark energy is required \textbf{[S]}:
	
	\begin{itemize}
		\item No accelerated expansion to explain
		\item Supernova observations explainable by wavelength-dependent redshift (original state; cf.\ K6 in Doc.~190)
		\item CMB anisotropies arise from $\xi$-field fluctuations, not primordial density perturbations
	\end{itemize}
	
	\section{Cosmic Verification through the CMB\_En.py Script}
	
	\subsection{Automated Calculations}
	
	The Python verification script \texttt{CMB\_En.py} (available on GitHub: ) performs calculations of the T0-cosmological relations named here:
	
	\begin{itemize}
		\item \textbf{Characteristic $\xi$-Length Scale:} $\Lxi = 100\,\mu\text{m}$
		\item \textbf{CMB-Temperature Verification:} Theoretical vs. experimental
		\item \textbf{Casimir-CMB Ratio:} 102.8 (or 103.8 with rounded $\Lxi = 100\,\mu$m; 1.0\,\% deviation from rounding, not a comparison with a measurement) \textit{(Corrected on 1~Oct.~2026: previously ``Agreement of 98.7\,\% (1.3\,\% deviation)'' -- see the section on the Casimir-CMB ratio.)}
		\item \textbf{Scaling Behavior:} Tested over 5 orders of magnitude
		\item \textbf{Energy Density Consistency:} Dimensional analysis of the energy densities
	\end{itemize}
	
	\begin{experiment}
		\textbf{Automated Verification of T0-Cosmology:}
		
		The script generates:
		\begin{itemize}
			\item Detailed log files with all calculation steps
			\item Markdown reports for scientific documentation
			\item LaTeX documents for publications
			\item JSON data export for further analyses
		\end{itemize}
		
		\textbf{Result:} The computed quantities lie close to the comparison values (Casimir-CMB ratio: 1.0\,\% deviation, solely from rounding $\Lxi$). \textit{(Corrected on 1~Oct.~2026: previously $1.3\,\%$ -- consequence of the correction $240 \to 720$.)}
	\end{experiment}
	
	\subsection{Reproducible Science}
	
	The automation of T0 calculations supports:
	
	\begin{itemize}
		\item \textbf{Transparency:} All calculation steps documented
		\item \textbf{Reproducibility:} Identical results on every run
		\item \textbf{Scalability:} Easy extension for new tests
		\item \textbf{Validation:} Automatic consistency checks
	\end{itemize}
	
	\section{Philosophical Implications}
	
	\subsection{The World View of T0-Cosmology}
	
	\begin{revolutionary}[Key statement]
		\textbf{Interpretation of T0-Cosmology [S]:}
		
		The document interprets the universe as the expression of a mathematical order that is described at its core by the parameter $\xipar$.
	\end{revolutionary}
	
	Possible philosophical conclusions from this interpretation:
	
	\begin{itemize}
		\item \textbf{Eternal Existence:} The universe had no beginning and will have no end
		\item \textbf{Mathematical Order:} All structures follow exact geometric principles
		\item \textbf{Universal Unity:} Quantum and cosmic scales are fundamentally connected
		\item \textbf{Deterministic Evolution:} Randomness is excluded at the fundamental level
	\end{itemize}
	
	\subsection{Epistemological Significance}
	
	The T0-Theory stands for the view that:
	
	\begin{itemize}
		\item Complex phenomena can be derived from simple principles
		\item Mathematical simplicity is a useful heuristic criterion
		\item Reductionism to a fundamental parameter is possible
		\item The universe is rationally comprehensible
	\end{itemize}

	\subsection{Technological Applications}
	
	Conceivable, purely speculative technological consequences \textbf{[S]}:
	
	\begin{itemize}
		\item \textbf{$\xi$-Field Manipulation:} Control over fundamental vacuum properties
		\item \textbf{Energy Extraction:} Tapping into the cosmic $\xi$-field
		\item \textbf{Communication:} $\xi$-based instantaneous information transfer
		\item \textbf{Transport:} $\xi$-field-supported propulsion systems
	\end{itemize}

	
	\begin{thebibliography}{30}
		
		\bibitem{t0_basics}
		Pascher, J. (2025). 
		\textit{T0-Theory: Fundamental Principles}. 
		T0 Document Series, Document 1.
		
		\bibitem{t0_gravitationalconstant}
		Pascher, J. (2025). 
		\textit{T0-Theory: Gravitational Constant}. 
		T0 Document Series, Document 3.
		
		\bibitem{t0_particlemasses}
		Pascher, J. (2025). 
		\textit{T0-Theory: Particle Masses}. 
		T0 Document Series, Document 4.
		
		\bibitem{cmb_verification_script}
		Pascher, J. (2025). 
		\textit{T0-Model Casimir-CMB Verification Script}. 
		GitHub Repository. 

		\bibitem{cosmic_document}
		Pascher, J. (2025). 
		\textit{T0-Theory: Cosmic Relations}. 
		Project Documentation. 

		\bibitem{heisenberg1927}
		Heisenberg, W. (1927). 
		\textit{On the Perceptual Content of Quantum Theoretical Kinematics and Mechanics}. 
		Zeitschrift für Physik, 43(3-4), 172--198.
		
		\bibitem{planck2020}
		Planck Collaboration (2020). 
		\textit{Planck 2018 results. VI. Cosmological parameters}. 
		Astronomy \& Astrophysics, 641, A6.
		
		\bibitem{casimir1948}
		Casimir, H. B. G. (1948). 
		\textit{On the attraction between two perfectly conducting plates}. 
		Proceedings of the Royal Netherlands Academy of Arts and Sciences, 51(7), 793--795.
		
		\bibitem{lamoreaux1997}
		Lamoreaux, S. K. (1997). 
		\textit{Demonstration of the Casimir force in the 0.6 to 6 $\mu$m range}. 
		Physical Review Letters, 78(1), 5--8.
		
		\bibitem{riess2022}
		Riess, A. G., et al. (2022). 
		\textit{A Comprehensive Measurement of the Local Value of the Hubble Constant}. 
		The Astrophysical Journal Letters, 934(1), L7.
		
		\bibitem{weinberg1989}
		Weinberg, S. (1989). 
		\textit{The cosmological constant problem}. 
		Reviews of Modern Physics, 61(1), 1--23.
		
		\bibitem{peebles2003}
		Peebles, P. J. E. (2003). 
		\textit{The Lambda-Cold Dark Matter cosmological model}. 
		Proceedings of the National Academy of Sciences, 100(8), 4421--4426.
		
		\bibitem{einstein1917}
		Einstein, A. (1917). 
		\textit{Cosmological Considerations on the General Theory of Relativity}. 
		Sitzungsberichte der Königlich Preußischen Akademie der Wissenschaften, 142--152.
		
		\bibitem{hubble1929}
		Hubble, E. (1929). 
		\textit{A relation between distance and radial velocity among extra-galactic nebulae}. 
		Proceedings of the National Academy of Sciences, 15(3), 168--173.
		
		\bibitem{friedmann1922}
		Friedmann, A. (1922). 
		\textit{On the Curvature of Space}. 
		Zeitschrift für Physik, 10(1), 377--386.
		
	\end{thebibliography}
	

